\documentclass[11pt,a4paper]{article}
\usepackage[T1]{fontenc}
\usepackage[utf8]{inputenc}
\usepackage[round,authoryear]{natbib}
\usepackage[french]{babel}
\usepackage[margin=25mm]{geometry}
\usepackage{lmodern}
\usepackage{microtype}
\usepackage{booktabs,longtable,array}
\usepackage[hidelinks,breaklinks]{hyperref}
\usepackage{url}
\setlength{\parindent}{0pt}
\setlength{\parskip}{6pt}
\setlength{\emergencystretch}{3em}
\clubpenalty=10000
\widowpenalty=10000
\renewcommand{\arraystretch}{1.2}
\newcommand{\titrethese}{Effet de l'incorporation des « extraits » de feuilles d'olivier dans l'eau de boisson du poulet chair}
\hypersetup{pdftitle={Dossier bibliographique initial - Mohamed Hamida},pdfauthor={Mohamed Hamida},pdfsubject={Projet de thèse vétérinaire 2026/2027}}

\begin{document}
\hypersetup{pageanchor=false}
\begin{titlepage}
\centering
{\large École nationale de médecine vétérinaire\\de Sidi Thabet -- Tunisie\par}
\vspace{1.4cm}
{\large Année universitaire 2026/2027\par}
\vspace{1.6cm}
{\LARGE\bfseries \titrethese\par}
\vspace{1.5cm}
{\Large Mohamed Hamida\par}
\vspace{1.5cm}
{\Large\bfseries Dossier bibliographique initial\par}
{\large Périmètre, premières références et plan de rédaction\par}
\vfill
\begin{minipage}{0.88\textwidth}
\small
Document de travail du 23 septembre 2026. Le titre reprend celui de la thèse PDF fournie. Le périmètre retenu couvre la préparation et la dose d'extrait, les polyphénols, la croissance, le gain moyen quotidien, l'indice de consommation, l'intégrité intestinale, le microbiote, les marqueurs sanguins et l'immunité. Le protocole et les résultats de Mohamed restent à fournir. Ce dossier constitue un premier repérage bibliographique.
\end{minipage}
\vspace{1cm}
\end{titlepage}
\hypersetup{pageanchor=true}

\section{Objet et méthode de l'actualisation}

La future thèse sera rédigée en français et reprendra l'organisation générale des deux sources ENMV : introduction, partie bibliographique, partie expérimentale, discussion, conclusion, références et annexes. La priorité actuelle est de construire une synthèse critique des travaux sur les feuilles d'olivier chez le poulet de chair, en donnant la première place à l'administration dans l'eau de boisson.

La thèse de Ben Brahim Amine, soutenue en 2018, aborde notamment la composition des feuilles, leurs propriétés, l'extraction, les performances zootechniques et certains paramètres sanguins. Celle de Heni Nabil, soutenue en 2013, développe aussi l'alimentation pratique du poulet. Ces documents servent de références locales et de modèles d'organisation ; le nouveau texte sera rédigé à partir de sources scientifiques vérifiées.

Le repérage initial porte sur les publications de 2018 au 23 septembre 2026, avec conservation d'études antérieures directement pertinentes. Il combine recherche web, pages d'éditeurs, notices PubMed, consultation de textes accessibles et validation des métadonnées par Crossref. Les interrogations structurées des bases bibliographiques, la recherche des citations et le contrôle des corrections ou rétractations restent à compléter. Une nouvelle actualisation sera nécessaire avant le dépôt en 2027.

Le périmètre retenu inclut désormais explicitement les polyphénols, le gain moyen quotidien (GMQ), l'indice de consommation (IC), l'intégrité intestinale et le microbiote. Les requêtes correspondantes sont préparées dans le protocole de recherche ; leur exécution et l'élargissement de la sélection restent à réaliser. Ce choix thématique ne préjuge pas des mesures effectivement réalisées dans l'essai de Mohamed.

Les études sur les extraits dans l'aliment, les feuilles entières, les poudres et les autres produits oléicoles seront classées séparément. Elles éclairent la question mais ne sont pas directement interchangeables avec une infusion administrée dans l'eau. Le volume d'extrait ajouté ne suffit pas, à lui seul, à établir une exposition comparable entre deux essais : sa composition et sa concentration doivent être documentées.

\section{Première lecture critique}

\subsection{Une priorité aux données obtenues dans l'eau}

Le travail de \citet{Jabri2017}, issu de l'ENMV de Sidi Thabet, constitue un point de départ local. Sa lecture nécessite toutefois des vérifications concernant la cohérence de l'effectif, les formulations des résultats et certaines indications de significativité. Les conclusions générales ne doivent donc pas remplacer l'examen des tableaux.

Une étude plus récente est particulièrement proche du sujet : \citet{Erener2023} ont évalué une infusion de feuilles dans l'eau de boisson. Selon le résumé vérifié, certains critères histomorphologiques intestinaux étaient améliorés, sans modification significative des performances de croissance ni des dénombrements cæcaux étudiés. Le texte intégral reste à obtenir pour l'évaluation complète. Cette distinction entre critères de réponse évite de qualifier globalement un extrait d'efficace ou d'inefficace.

\subsection{Croissance, efficacité alimentaire et marqueurs sanguins}

Les études dans l'aliment apportent des éléments complémentaires. \citet{Erener2020} ont étudié la croissance et des paramètres lipidiques, mais les modifications sanguines ne peuvent pas être décrites comme uniformément favorables. Une incohérence de formulation concernant le LDL devra être résolue à partir du tableau correspondant avant toute citation quantitative.

Le résumé de \citet{Xie2022} rapporte des effets sur les paramètres antioxydants et la viande, ainsi qu'une diminution du gain et de l'ingestion pour un niveau d'incorporation. L'étude souligne l'intérêt d'une analyse par dose et par critère. La lecture intégrale est encore nécessaire. Les résultats neutres de \citet{Leskovec2018} sur l'utilisation des nutriments rappellent également que les bénéfices ne sont pas systématiques.

\subsection{Statut oxydatif et immunité}

Les effets antioxydants ne constituent pas, à eux seuls, une preuve d'amélioration de l'immunité. \citet{Pirman2021} rapportent une diminution du MDA plasmatique dans un contexte alimentaire particulier, sans amélioration significative de la croissance. Des divergences entre résumé et tableaux, notamment sur le 8-OHdG et les acides gras à chaîne courte, justifient une lecture attentive des résultats. Les études sur la conservation de la viande, telles que \citet{Sarica2018}, doivent être situées dans leur domaine d'application.

Le travail de \citet{Mahasneh2024} comporte des mesures sanguines et immunitaires dans un contexte de stress thermique. Il utilise des feuilles incorporées à l'aliment et plusieurs groupes de plantes. Il faudra distinguer les effets du groupe olivier de ceux des associations, puis examiner si le contexte correspond à l'essai de Mohamed. Les modalités réelles de ses mesures immunitaires ne sont pas encore connues.

\citet{Alfifi2025} associent l'extrait phénolique de feuilles à l'arginine et étudient notamment des critères immunitaires. L'absence de groupe recevant l'extrait seul empêche d'attribuer directement à celui-ci l'ensemble des effets observés. Cette étude est une piste prioritaire pour la lecture approfondie de l'axe immunité.

\subsection{Caractérisation et nouveautés}

Les travaux de \citet{Vasilopoulou2023} et \citet{Vasilopoulou2026} concernent respectivement un extrait purifié et un extrait enrichi en oléuropéine dans l'aliment. Ils invitent à documenter la nature du produit étudié et à vérifier l'indépendance des essais publiés par une même équipe. La référence de 2026 possède un DOI contenant 2025 : l'année du volume, la date de mise en ligne et le millésime du DOI doivent être distingués.

Une publication avancée du 19 août 2026, \citet{Hiba2026}, compare un extrait libre à un extrait encapsulé dans l'aliment. Son résumé et ses métadonnées ont été vérifiés. Elle intéresse directement la préparation et les marqueurs sanguins ; une éventuelle amélioration de la biodisponibilité ne doit pas être affirmée avant vérification du texte intégral et des mesures réalisées.

\clearpage
\section{Plan proposé pour la partie bibliographique}

\begin{enumerate}
\item L'olivier et la valorisation de ses feuilles en Tunisie : contexte concis et données récentes vérifiées.
\item Polyphénols des feuilles : composition, variabilité et méthodes de caractérisation.
\item Préparation, caractérisation, stabilité et administration des extraits ; comparabilité des doses.
\item Performances zootechniques : croissance, GMQ, ingestion et IC ; périodes de calcul, unités et correction éventuelle de la mortalité.
\item Intégrité intestinale : histomorphologie, marqueurs de barrière et mesures fonctionnelles, selon les méthodes des études.
\item Microbiote intestinal et cæcal : dénombrements ciblés, profils de communautés et méthodes d'analyse.
\item Marqueurs sanguins et statut oxydatif : résultats, sens biologique et limites.
\item Immunité : nature des critères étudiés et niveau de preuve.
\item Synthèse critique, divergences entre études et question de recherche propre à la thèse.
\end{enumerate}

La synthèse distinguera les dénombrements bactériens par culture des analyses de communautés par séquençage, et les critères morphologiques des mesures fonctionnelles de la barrière intestinale. Les méthodes, sites et dates de prélèvement seront conservés pour comparer les études. Les résultats relatifs aux métabolites microbiens seront présentés séparément de ceux décrivant la composition du microbiote.

Le protocole réel permettra de préciser les hypothèses et les critères expérimentaux. Les doses, les groupes, les répétitions, les dates de mesure et les résultats ne seront renseignés qu'à partir des documents de Mohamed. Les axes retenus guident dès maintenant la recherche et la rédaction bibliographiques.

\clearpage
\section{État de consultation des références}

Le registre détaillé, conservé dans \texttt{bibliography/reading\_register.json}, distingue lecture intégrale, consultation partielle et résumé seul. La présence d'une référence ci-dessous ne signifie pas que son évaluation critique est achevée. Le fichier \texttt{references.bib} conserve les notices ; le manifeste des téléchargements indique l'origine des PDF et leur empreinte numérique.

\begin{longtable}{@{}>{\raggedright\arraybackslash}p{3.4cm}>{\raggedright\arraybackslash}p{3.2cm}>{\raggedright\arraybackslash}p{8.3cm}@{}}
\toprule
Référence & Administration & Usage et contrôle restant \\
\midrule
\endhead
Jabri et al., 2017 & Eau & Point de départ local ; incohérences de présentation à vérifier. \\
Leskovec et al., 2018 & Aliment & Nutriments et os ; résultats neutres et contexte alimentaire. \\
Sarıca et Ürkmez, 2018 & Aliment & Qualité de viande ; ne démontre pas une réponse immunitaire. \\
Erener et al., 2020 & Aliment & Croissance et lipides ; vérifier le sens de variation du LDL. \\
Pirman et al., 2021 & Aliment & Statut oxydatif ; vérifier les résultats dans les tableaux. \\
Xie et al., 2022 & Aliment & Dose et critères de réponse ; texte intégral à obtenir. \\
Erener et al., 2023 & Eau & Étude prioritaire ; résumé vérifié, texte intégral à obtenir. \\
Vasilopoulou et al., 2023 & Aliment & Extrait purifié ; lecture intégrale à compléter. \\
Mahasneh et al., 2024 & Aliment & Feuilles et stress thermique ; portée contextuelle. \\
Alfifi et al., 2025 & Aliment & Association à l'arginine ; absence de groupe extrait seul. \\
Vasilopoulou et al., 2026 & Aliment & Extrait enrichi ; texte intégral et indépendance de l'essai à vérifier. \\
Hıba et al., 2026 & Aliment & Encapsulation ; publication avancée, texte intégral à obtenir. \\
\bottomrule
\end{longtable}

\clearpage
\renewcommand{\bibfont}{\small}
\setlength{\bibsep}{3pt}
\bibliographystyle{abbrvnat}
\bibliography{../bibliography/references}
\end{document}
